\section{Final checklist}

This chapter developed conditional probability as probability after information
has changed the active space of reasoning. The central habit is to identify the
reference space before calculating.

The checklist below summarizes the method.

\subsection*{Step 1: Identify the original model}

Before conditioning, there must be a probability model. Identify:
\begin{itemize}
    \item the sample space \(\Omega\);
    \item the elementary outcomes or relevant events;
    \item the probability assignment;
    \item whether the model is uniform or non-uniform.
\end{itemize}

Conditional probability does not replace model construction. It uses the model
that has already been built.

\subsection*{Step 2: Identify the information event}

In an expression
\[
P(A\mid B),
\]
the event after the vertical bar is the information event. It is the event that
is assumed to have occurred.

Say explicitly:
\[
\text{we are now reasoning inside }B.
\]

Check that
\[
P(B)>0.
\]
If \(P(B)=0\), the elementary conditional probability formula is not defined.

\subsection*{Step 3: Identify the target event}

The event before the vertical bar is the event being measured. In
\[
P(A\mid B),
\]
the target event is \(A\). The favourable part of the information region is
\[
A\cap B.
\]

Do not replace \(A\) by \(B\). Do not reverse the order unless the problem asks
for the reverse conditional probability.

\subsection*{Step 4: Use the definition when conditioning directly}

If the problem asks directly for \(P(A\mid B)\), use
\[
P(A\mid B)=\frac{P(A\cap B)}{P(B)}.
\]
In a finite uniform model this becomes
\[
P(A\mid B)=\frac{|A\cap B|}{|B|}.
\]
The denominator is the size or probability of the information region. The
numerator is the part of that region where the target event also occurs.

\subsection*{Step 5: Use independence only when justified}

Events \(A\) and \(B\) are independent when information about one does not change
the probability of the other. Algebraically,
\[
P(A\cap B)=P(A)P(B).
\]
When probabilities are positive, this is equivalent to
\[
P(A\mid B)=P(A)
\]
and
\[
P(B\mid A)=P(B).
\]

Do not assume independence merely because there are two stages or two
coordinates. Independence belongs to the probability assignment, not just to the
notation.

\subsection*{Step 6: Use total probability when cases produce the event}

If an event \(A\) may occur through several mutually exclusive and exhaustive
cases \(B_1,
\ldots,B_n\), use the law of total probability:
\[
P(A)=\sum_{i=1}^n P(A\mid B_i)P(B_i).
\]
Before using it, check:
\begin{itemize}
    \item the cases are disjoint;
    \item the cases cover the whole sample space;
    \item the conditional probabilities inside the cases are meaningful.
\end{itemize}

Read the formula as:
\[
\text{total probability}=\sum \text{case probability}\times\text{probability inside case}.
\]

\subsection*{Step 7: Use Bayes' theorem when reversing information}

Use Bayes' theorem when the problem gives or suggests probabilities of the form
\[
P(\text{observation}\mid\text{case})
\]
but asks for
\[
P(\text{case}\mid\text{observation}).
\]

For a partition \(B_1,
\ldots,B_n\),
\[
P(B_j\mid A)=
\frac{P(A\mid B_j)P(B_j)}
{\sum_{i=1}^n P(A\mid B_i)P(B_i)}.
\]

Remember the roles:
\begin{itemize}
    \item \(P(B_j)\) is the prior probability of the case;
    \item \(P(A\mid B_j)\) is the likelihood of the observation in that case;
    \item \(P(B_j\mid A)\) is the posterior probability of the case after the
    observation.
\end{itemize}

\subsection*{Step 8: Interpret the result}

After computing a conditional probability, translate it back into words.

For example,
\[
P(A\mid B)=0.7
\]
means:
\[
\text{among the outcomes compatible with }B,\text{ the event }A\text{ has probability }0.7.
\]
It does not mean that \(B\) was caused by \(A\), nor that \(P(B\mid A)=0.7\).

\subsection*{Compact checklist}

For any problem in this chapter, ask:

\begin{enumerate}
    \item What is the original probability model?
    \item What information has been given?
    \item What event is after the vertical bar?
    \item What event is before the vertical bar?
    \item What is the intersection of the target event with the information
    event?
    \item Is the conditioning event positive-probability?
    \item Is independence stated or justified?
    \item Is there a partition of cases?
    \item Is the problem asking for a direct conditional probability or a reverse
    probability requiring Bayes' theorem?
    \item What does the final number mean in words?
\end{enumerate}

\begin{theorembox}
The most important sentence in conditional probability is:
\[
\text{the event after the vertical bar is the space in which we now reason.}
\]
Once this is clear, the formulas become expressions of that change of reference
space.
\end{theorembox}
